\begin{abstract}
When measurement errors are correlated, selecting observations changes the
covariance matrix that must be inverted: the information of a schedule $S$
is $J(S)=J_0+F_S^\top R_{SS}^{-1}F_S$, generally not a sum of per-observation
terms. We
develop global certificates for this discrete design problem with Gaussian
errors, known covariance and nominal mean sensitivities. First, a common
surrogate built from entries of the full precision matrix overstates
information by at most the Kantorovich factor of the covariance, which gives a
sharp efficiency bound; in a public kinetics model it changes one of 33
exactly enumerated criterion--budget optima. Second, finite-history Gaussian
conditionals satisfy explicit relative precision bounds, uniform over
schedules, for noisy scalar chains, complete observed blocks, exponentially
decaying covariances and fixed partial-observation packets. On a
calendar-state graph they give a weighted-trace FPTAS with input-sized
dimensions when the promise constants are fixed. For fixed information dimension and
additive rational PSD atoms, polynomial sets of feasible paths in explicit
acyclic graphs, or of bases of rationally represented matroids, approximate
every feasible information matrix within $1\pm\eta$ in PSD order with exact
ranges, which yields D-, A-, E- and contrast approximation schemes. Third, we
give rationally checked certificates: tangent support bounds, one witness
that bounds every admissible diagonal virtual-noise relaxation from below,
certificates for one schedule over finite scenarios with unknown optimal
normalizers, and a nested latent-anchor hierarchy proved by Schur
elimination. On ten stipulated rational instances, finite-history upper bounds lie
strictly below the optimum of every admissible scalar virtual-noise
relaxation, and on one instance below that of every diagonal one.
Full-packet and fixed-physical-grid experiments show limits in bound
strength and cost.
\end{abstract}

\noindent\textbf{Keywords:} optimal experimental design; correlated errors;
discrete optimization; relative matrix approximation; global certificates;
latent Gaussian models.

\section{Introduction}\label{sec:introduction}

Measurement selection combines a statistical decision with a discrete resource
allocation problem. A kinetic experiment may choose which species to measure,
when to sample and which instruments to install; a sensor schedule may acquire
a whole multichannel packet whenever a device is activated. Installation
costs, exact sample counts and minimum time gaps restrict the feasible
experiments. Correlated errors couple these choices: the information from a
new observation depends on which other observations are acquired.
Correlated dynamic multiresponse design has a long history \citep{Patan2007},
and recent process-design formulations integrate measurement costs with
model-based information criteria \citep{Wang2024}. For optimization, the
central question is how to bound the best attainable information over all
feasible discrete experiments.

Write $F_S$ for the nominal mean sensitivities of the selected observations,
$R_{SS}$ for their known positive definite error covariance, and $J_0$ for
fixed prior information or a declared regularizer. The information matrix is
$J(S)=J_0+F_S^\top R_{SS}^{-1}F_S$. Even with a linear mean, $J(S)$ is generally not a sum of fixed per-observation terms.
Taking selected entries of the full precision matrix $R^{-1}$ instead gives a
different information model, as \citet{Liu2016} already explained. A valid
optimization bound must refer to the experiment that will actually be
performed.

Even with the correct information formula, two further gaps must be kept
separate from the numerical optimization gap. A continuous mixture of
feasible schedules can have a larger log determinant than any single
schedule, and closing the numerical gap does not close this integrality gap.
Approximating the covariance adds an error that must be controlled uniformly
over all schedules the optimizer may choose. In this paper a
\emph{certificate} is a feasible design with a checked lower bound on its
objective, together with a global upper bound for the same discrete model.
For $p$ parameters, a log-determinant gap $g$ certifies determinant-root
efficiency at least $e^{-g/p}$, whether the design came from a heuristic, a
convex relaxation or an exact combinatorial search.

\subsection{Relationship to prior work}\label{subsec:intro-prior}

\paragraph{Correlated design and sensor selection.}
The selected-covariance information $J(S)$ is the established model in
correlated design and sensor selection \citep{Patan2007,Liu2016,Pazman2022}.
\citet{Liu2016} also analyze the weak-correlation regime of gating. Virtual-noise designs and covariance-split sensor
relaxations give concave continuous extensions \citep{Liu2016,Pazman2022};
\citet[Proposition 3]{Hainy2025} prove that Liu's scalar split and the
virtual-noise design coincide, and optimize them by simplicial decomposition. Information-matrix hulls
\citep{HarmanTrnovska2009}, subsystem criteria and mixed-integer
second-order cone formulations for exact designs \citep{Sagnol2015}, and
exact multiresponse designs \citep{Filova2026} are also established. We use
these results as building blocks and comparators; our comparisons concern
the strength of specified relaxation families, not dense formulations in
general.

\paragraph{Local Gaussian conditionals.}
Local conditional factorizations originate with \citet{Vecchia1988}. Their
noisy and vector extensions, sparse inverse-factor interpretations and
approximation theory are substantial existing work
\citep{KatzfussGuinness2021,Schafer2021,Huan2025,KaminetzWebber2026}, and
\citet{Kozdoba2019} use process-noise dissipation to obtain finite-memory
Kalman forecasts. If the error process is fully observed and Markov, exact
gap conditionals give an ordered path representation whose inverse-matrix
hull is covered by the convexification results of \citet{Lee2026}, building
on \citet{Wei2023}. Independent measurement noise generally destroys this
observed Markov property. We treat the noisy case through the actual local
regressions, with explicit constants for chronological calendar windows; we
do not introduce a new factorization or a general localization principle.

\paragraph{Approximation algorithms.}
Bounded-width Gaussian message schemes \citep{Mahalanabis2012}, approximate
Pareto sets \citep{Papadimitriou2000,Tsaggouris2009,Mittal2013}, nonlinear
matroid optimization \citep{Berstein2008} and fixed-dimensional design
algorithms \citep{Brown2024} are the main algorithmic precedents. A scalar
focused-target reduction of the Gaussian-message scheme already yields an
approximation scheme (Appendix~\ref{app:scalar-prior}), so we do not claim a
first scalar FPTAS. Greedy guarantees for weakly supermodular objectives
achieve sparse multiple-regression bounds with an enlarged support
\citep{Liberty2017}; our schemes enforce the exact cardinality and do not
rely on sub- or supermodularity. Appendix~\ref{app:pivot-supermodularity}
shows that a supermodularity step asserted for a pivot-set KL objective in
\citet{KaminetzThesis2025} fails. Spectral spanners and coresets
\citep{Indyk2020,Mahabadi2026} retain elements rather than returning one
feasible object close to each target matrix. For A- and E-design with
additive rank-one information under partition constraints,
\citet{BansalXu2026} rule out polynomial-factor approximation in a growing
information dimension unless $\mathsf{P}=\mathsf{NP}$, consistent with our
fixed-dimensional results.

\paragraph{Robustness, nuisance parameters and certification.}
Efficiency normalization over scenarios \citep{Burclova2016,Maus2010},
robust and scenario-based design
\citep{Duarte2018,WangYue2026,Chowdhary2025}, and nuisance or subsystem
design \citep{Alexanderian2021,Sagnol2015,LevineHow2013}
precede our robust and latent-anchor certificates. Log-determinant support
conditions \citep{HarmanTrnovska2009} and convexity of
$\log\det(I+KX^{-1})$ \citep{KimKim2006,Kim2019} supply the tangent and dual
inequalities we check in rational arithmetic. Disjunctive programming
explains why grouping alternatives before convexification strengthens a
relaxation \citep{Balas1985,Papageorgiou2025}.

Table~\ref{tab:prior-contributions} lists what is inherited and what is
developed here; detailed comparisons accompany the relevant proofs.

\begin{table}[htbp]
\centering\small
\caption{Established antecedents and the developments in this paper. The new
results combine established statistical and optimization tools under the
hypotheses stated in the cited sections.}
\label{tab:prior-contributions}
\begin{tabular}{@{}>{\raggedright\arraybackslash}p{.20\textwidth}>{\raggedright\arraybackslash}p{.34\textwidth}>{\raggedright\arraybackslash}p{.40\textwidth}@{}}
\toprule
Component & Established antecedent & Development in this paper\\
\midrule
Selected information & Selected covariance and gating distinction
\citep{Liu2016}; classical Kantorovich inequality \citep{Moradi2019}
& Sharp design-efficiency consequences, equality/rank analysis, and exact
public-model reanalysis (Sections~\ref{sec:foundations},
\ref{subsec:source-computation}).\\[4pt]
Finite history & Vecchia conditionals, inverse localization and covariance
transport \citep{Vecchia1988,Schafer2021,Kozdoba2019}
& Explicit schedule-uniform relative bounds for four covariance classes,
spacing refinements, and transfer to exact feasible support prices
(Sections~\ref{sec:locality}, \ref{subsec:local-certificates}).\\[4pt]
Spectral approximation & Pareto/profile algorithms and guessed-vector
normalization \citep{Berstein2008,Tsaggouris2009,Brown2024}
& Deterministic all-target, two-sided PSD covers with polynomial accuracy
dependence and exact ranges for explicit rational acyclic graphs and represented matroids
(Section~\ref{sec:approximation}).\\[4pt]
Global comparison & Virtual-noise concavity and equivalence
\citep{Liu2016,Pazman2022,Hainy2025}; inverse-logdet convexity
\citep{KimKim2006}
& Finite rational witnesses separating every scalar or diagonal split,
including their stated affine-support intersections
(Section~\ref{subsec:all-split}).\\[4pt]
Robustness and anchors & Efficiency normalization
\citep{Burclova2016,Maus2010}; nuisance and subsystem design
\citep{Alexanderian2021,Sagnol2015}
& Certified uncertainty in optimal normalizers for one common schedule, and a proved
ordering across nested changing-anchor representations
(Sections~\ref{subsec:robust-certificate}, \ref{subsec:separators}).\\
\bottomrule
\end{tabular}
\end{table}

\subsection{Main results}\label{subsec:intro-results}

\begin{enumerate}
\item \emph{Selected versus gated information.} Full-inverse gating never
understates information. If $mI\preceq R\preceq MI$ and $\kappa=M/m$,
then $J(S)\preceq J_{\gate}(S)\preceq\alpha J(S)$ with the Kantorovich
factor $\alpha=(\kappa+1)^2/(4\kappa)$
(Propositions~\ref{prop:gating} and~\ref{prop:kantorovich}). A global
gated D-optimum therefore has determinant-root efficiency at least
$1/\alpha$, and this bound is sharp as an infimum, already for one
parameter (Corollary~\ref{cor:gated-design},
Example~\ref{ex:gated-sharpness}). In the public kinetics implementation of
\citet{Wang2024}, exact evaluation of all 2347 schedules changes one of 33
criterion--budget optima: the trace-information choice at one budget, with
a relative loss of about $3.11\%$ (Section~\ref{subsec:source-computation}).

\item \emph{Schedule-uniform locality bounds.} If the normalized residual
covariance of the local conditionals is within $\delta$ of the identity, the
local precision lies between $1-\delta$ and $1+\delta$ times the true
precision (Lemma~\ref{lem:residual-transfer}). We bound $\delta$
explicitly, for every selected schedule at once, for noisy scalar chains
(Theorem~\ref{thm:scalar-locality}), complete observed blocks with
noncommuting transitions (Theorem~\ref{thm:block-locality}), general
covariances with exponential off-diagonal decay and a spectral floor
(Theorem~\ref{thm:general-decay}), and fixed partial-observation packets,
including singular latent covariances (Theorem~\ref{thm:partial-locality}).
The bounds decay geometrically in the window length and contain no
dimension factor. For stationary scalar covariances,
Proposition~\ref{prop:spacing-locality} refines them under minimum
sampling gaps.

\item \emph{Approximation with exact feasibility.} On a calendar-state graph,
these bounds give a rational weighted-trace FPTAS with input-sized packet
and parameter dimensions under fixed promises, including exact cardinality,
mandatory and forbidden times, and minimum spacing
(Theorem~\ref{thm:trace-fptas}). Selecting individual channels within a
packet admits no deterministic FPTAS for scalar information unless
$\mathrm P=\mathrm{NP}$, even at a single time with a well-conditioned
covariance (Proposition~\ref{prop:channel-hardness}). For fixed information
dimension, Theorems~\ref{thm:dag-spectral-set} and~\ref{thm:matroid-spectral-set}
return polynomial sets of actual paths in explicit acyclic graphs, or of
bases of rationally represented matroids, with additive rational PSD atoms.
For every feasible matrix, one returned matrix lies between its $1-\eta$ and
$1+\eta$ multiples in PSD order and has the same range, even when singular.
The proof combines rational normalization on exact ranges with bounded
signed profiles and exact feasible-witness recovery. The covers give
approximation schemes for D-, E- and conventional A-optimality and for
estimable contrasts, with no condition-number bound. Under the fixed
promises of Theorem~\ref{thm:trace-fptas}, the DAG cover transfers to the
true selected-covariance model (Corollary~\ref{cor:true-spectral-cover}).
Relative to \citet{Brown2024}, who give a randomized PTAS in the broader
independence-oracle matroid model, our represented-matroid result is
deterministic with polynomial dependence on $1/\eta$ and does not subsume
their model. To the best of our knowledge, this combined guarantee is not
stated in the inspected predecessors. The polynomial degree grows with the
information dimension.

\item \emph{Certificates and relaxation comparisons.} Any positive definite
reference matrix and a globally priced linear support over the feasible
graph give an upper bound on the original discrete objective
(Theorem~\ref{thm:local-certificate}); rational arithmetic checks the
reference, the full support and outward logarithm bounds. A fixed fractional
point and a repaired PSD dual give one lower witness valid for every
admissible diagonal virtual-noise split
(Proposition~\ref{prop:diagonal-barrier}).
A strict separation from a valid integer upper bound therefore survives both
complete continuous optimization and split tuning within that family. For
one schedule shared by finitely many scenarios, we certify standardized
D-efficiency while enclosing each scenario's unknown optimal normalizer
(Proposition~\ref{prop:robust-normalizers}). Under the fixed memory
promises and feasibility conditions of Theorem~\ref{thm:trace-fptas}, with
fixed scenario count and information dimension, this problem has an FPTAS
(Corollary~\ref{cor:robust-fptas}).

\item \emph{Latent separators.} When long calendar histories are expensive,
latent anchors give an exact representation by conditional blocks, with its
own support certificate (Theorem~\ref{thm:separator-certificate}). For
nested anchor sets, removing anchors never loosens the complete-schedule
hull bound (Theorem~\ref{thm:separator-hierarchy}), by a cross-representation
Schur identity. Within established Schur and subsystem-design theory
\citep{Sagnol2015}, the claimed novelty is this comparison across changing
latent representations, which to the best of our knowledge is not stated in
the inspected predecessors. Larger blocks cost more to enumerate, and even the
empty-anchor hull can retain an integrality gap.

\item \emph{Computational evidence} (Section~\ref{sec:computation}). On ten
rational instances (six synthetic, four kinetic), certified finite-history
log gaps are at most $0.0039$, and every admissible scalar virtual-noise
relaxation has optimum at least $0.022$ above the finite-history upper
bound. On one kinetic instance a single witness separates every admissible
diagonal split by more than $0.092$. A three-scenario kinetic certificate
guarantees at least $99.6\%$ of the unknown optimal standardized robust
efficiency. An exact eight-time study verifies strict tightening across four
nested anchor sets, and a new far-pair argument improves the complete-block
constant. The limits are reported as well: in full-packet examples the dense
bound is tighter; on the finer fixed physical grids (96 and 192 points),
the primary calendar runs complete no pricing step within the 30-second allowance; and the
separator certificates miss the $0.01$ log-gap target.
\end{enumerate}

\paragraph{Scope.}
Throughout, the covariance is known and parameter independent, the
sensitivities are local, and a certificate concerns the supplied finite
experiment. PSD priors are allowed in the locality and spectral theory; the
practical log-determinant and separator certificates assume a positive
definite prior. The guarantees do not imply global nonlinear
identifiability or robustness outside a declared scenario set.

\paragraph{Formal verification.}
Lean formalizes the explicit-DAG and represented-matroid covers
(Theorems~\ref{thm:dag-spectral-set} and~\ref{thm:matroid-spectral-set})
from rational inputs, including exact singular ranges, criterion selection
and bit-work bounds; the DAG core assumes verified topological vertex
indices. The transfer in Corollary~\ref{cor:true-spectral-cover} is
formalized conditional on its sandwich hypothesis; the stochastic locality
estimates and the graph construction are not formalized
(Sections~\ref{subsec:spectral-consequences} and~\ref{subsec:matroid-cover}).

\paragraph{Organization.}
Sections~\ref{sec:foundations}--\ref{sec:computation} follow the order of
these results, and Section~\ref{sec:discussion} discusses scope; a
standalone supplement replays the certificates without a commercial solver.
